% Exact characteristic-two plane sextic case; proof for independent audit.
\subsection{The exact group for ordinary plane sextics in characteristic two}
\label{subsec:sextic-exact}

\begin{theorem}\label{thm:sextic-full-monodromy}
For the universal ordinary smooth plane sextic in characteristic two,
the geometric and arithmetic monodromy images on
$V=H^1_{\mathrm{et}}(C_{\bar\eta},\mathbf F_2)$ are both
$\mathrm{GL}_{10}(\mathbf F_2)$.  Thus $V$ is irreducible and its
determinant character is trivial.
\end{theorem}

\begin{proof}
The generic genus is ten, and ordinarity gives $2^{10}=1024$
geometric theta characteristics.  The canonical characteristic $B_C$
is the unique non-effective one.  It is invariant under the
geometric Galois group, also when a purely inseparable extension is
needed to define it as a line bundle.  Identifying the set of theta
characteristics with $V^\vee$ using $B_C$ as origin, the remaining
1023 effective theta characteristics correspond to the nonzero
vectors.  We first prove that monodromy is transitive on them.

\paragraph{The contact reciprocity calculation.}
Let $C$ be a smooth plane curve of degree $d$, let $Q$ be a smooth
plane curve of degree $e=d-3$, and suppose $C\cap Q=2D$.
The half-intersection divisor defines the same zero-dimensional
subscheme $D$ on both curves: locally, in coordinates in which one
curve is $y=0$, the other has contact $y=u x^{2m}+\cdots$, and both
half divisors have ideal $(y,x^m)$.
Put
\[
 \theta_C=\mathcal O_C(D),\qquad
 \theta_Q=\mathcal O_Q(D-3H).
\]
These are theta characteristics.  The two ideal-sheaf sequences,
twisted by $e$, are
\[
 0\to\mathcal O_{\mathbf P^2}(-3)\to\mathcal I_D(e)
       \to\theta_C\to0,
 \qquad
 0\to\mathcal O_{\mathbf P^2}\to\mathcal I_D(e)
       \to\theta_Q\to0.
\]
Since $H^1(\mathbf P^2,\mathcal O(-3))=H^1(\mathbf P^2,\mathcal O)=0$,
they give the exact numerical relation
\begin{equation}\label{eq:contact-reciprocity}
 h^0(\theta_C)=1+h^0(\theta_Q).
\end{equation}

\paragraph{Smooth contact cubics.}
For $d=6$, every effective $\theta_C$ with section $s$ supplies a
contact cubic, because $s^2\in H^0(C,\omega_C)=H^0(C,\mathcal O_C(3))$
lifts uniquely to a ternary cubic.
If $Q$ is smooth, its theta characteristics are degree-zero line
bundles.  For an ordinary cubic, the only non-effective one is $B_Q$;
the other is $\mathcal O_Q$.  For a supersingular cubic, the sole
geometric theta characteristic is $\mathcal O_Q$.

The incidence corresponding to $\theta_Q=B_Q$ over the ordinary
smooth-cubic parameter space is irreducible of dimension 27.
To see this without a descent assumption on $B_Q$, pull the
coefficient space back along Frobenius (and work in an affine
coefficient chart).  The quotient of the relative Frobenius pushforward
by the structure sheaf supplies $B_Q$ on the Frobenius-twisted universal
curve.  This purely inseparable base change is irreducible and changes
no geometric Galois group.  The spaces
$H^0(Q,B_Q(3))$ form a vector bundle of rank nine, since its degree is
nine and the genus is one.  Squaring their sections gives the allowed
restrictions in $H^0(Q,\mathcal O_Q(6))$.
The restriction of ternary sextics is surjective, with kernel
$Q\cdot H^0(\mathbf P^2,\mathcal O(3))$ of dimension ten.
Thus the allowed sextics form, after this purely inseparable cover,
a nineteen-dimensional affine space, or a projective space of
dimension eighteen.  Smooth cubics have parameter dimension nine,
so the total is $9+18=27$, equal to the sextic parameter dimension.
The image of an irreducible space is irreducible.

The smooth-contact-cubic loci with effective $\theta_Q$ do not
dominate the sextic parameter space.  Their incidence dimensions are
at most 27 by the same degree-nine Riemann--Roch calculation; by
\eqref{eq:contact-reciprocity}, every such fibre over a sextic contains
the positive dimensional projective linear system of its
$\theta_C$, which has at least two sections.  Their images therefore
have dimension at most 26.  This remains true for the supersingular
cubic locus, whose base dimension is smaller.

\paragraph{Singular contact cubics.}
We show that these also occur only for sextics in a proper closed
subset of the 27-dimensional parameter space.
For every fixed cubic, the kernel of restriction of sextic forms has
dimension ten.  The following bounds are therefore bounds on the
dimension of the allowed restriction space, before adding ten and
projectivizing.

An irreducible nodal cubic has normalization $\mathbf P^1$ and two
points $P_+,P_-$ above the node.  Its stratum has dimension eight.
The pullback of $\mathcal O_Q(6)$ has degree eighteen.
Away from $P_+,P_-$, the orders of the restricted sextic are even.
At these two points the parities are equal.  If both are even, the
binary form is a square of a degree-nine form; the node gluing
condition gives one independent linear condition on that square root,
leaving dimension nine.  If both are odd, it is a fixed binary
quadratic with zeros $P_+,P_-$ times a square of degree eight, again
dimension nine; the gluing values are then both zero.  The incidence
dimension is consequently at most $8+9+10-1=26$.

An irreducible cuspidal cubic has normalization $\mathbf P^1$, one
point above its cusp, and parameter dimension at most seven.  All
orders of the pulled-back sextic are even, since its total degree is
eighteen and all orders away from that point are even.  It is therefore
a square of a degree-nine binary form, a space of dimension ten.
Any descent condition can only decrease the dimension.  The resulting
bound is $7+10+10-1=26$.  The dimensions eight and at most seven of
these strata follow directly by fixing a singular point: singularity
imposes three independent linear conditions on cubic forms, and a
repeated tangent cone imposes at least one further condition.

If the cubic is a transverse union of a line and a smooth conic, its
parameter dimension is seven.  On the two normalization components,
the sextic has degrees six and twelve.  Its odd orders, if any, must
occur at both of their two intersection points.  With no odd orders,
the square-root spaces have dimensions four and seven, subject to
two independent gluing conditions, giving nine.  With odd orders at
both points, extracting their fixed divisors leaves square-root
spaces of dimensions three and six, also giving nine.  The total
incidence dimension is at most 25.  If the line is tangent to the
conic, the cubic stratum has dimension at most six; each component's
restriction must have even order at the unique intersection point,
and the two square-root spaces satisfy at least their common-value
condition.  Their combined dimension is at most ten, again giving
at most 25.

For three distinct lines in a triangle, the parameter dimension is
six.  The possible odd-node patterns are either no nodes or all three
nodes.  In the first case the three four-dimensional square-root
spaces have three independent common-value conditions; in the second,
the three spaces have dimension three each.  Both bounds are nine,
giving total dimension at most 24.  For three concurrent lines the
parameter dimension is five.  Each line restriction has even order
at their sole common point, and the three square-root spaces satisfy
two common-value conditions.  Their combined dimension is at most
ten, again giving at most 24.

Finally a nonreduced cubic is $L^2M$, allowing $L=M$.
Its restriction to a smooth sextic can be even only if the line
section $M|C$ is even.  Sextics with such a line lie in dimension at
most 26: for a fixed line, binary sextics that are squares form a
four-dimensional subspace of a seven-dimensional restriction space;
the ternary restriction kernel has dimension 21.  Projectivizing and
allowing the line to vary gives $(4+21-1)+2=26$.
All possible singular cubics have now been covered.

Outside the proper loci just bounded, every one of the 1023 effective
theta characteristics of an ordinary sextic has a smooth contact
cubic, and this cubic corresponds to $B_Q$ by
\eqref{eq:contact-reciprocity}.  The irreducible incidence therefore
dominates the sextic space.  Its dimension is 27, so its generic
fibre is finite; every effective theta has $h^0=1$, and the fibre has
exactly 1023 geometric points.  As in the quintic case, irreducibility
implies geometric Galois transitivity on these points even if the
incidence has a purely inseparable part.  Thus $G_{\mathrm{geom}}$
is transitive on $V^\vee\setminus\{0\}$.

\paragraph{A boundary element that excludes the smaller transitive groups.}
Degenerate a sextic to an ordinary smooth quintic and a transverse
line, meeting in five distinct points.  This is an ordinary nodal
sextic.  Its rank-ten local system has a four-dimensional graph
submodule and the rank-six quintic local system as quotient.
The node permutation group is a subgroup of $S_5$; the graph
representation is its augmentation representation.
The exact quintic monodromy theorem supplies a quotient with order
divisible by seven and an element of order thirty-one acting on the
six-dimensional quotient as a one-dimensional identity plus an
irreducible five-dimensional summand.  Hence $7\mid|G_{\mathrm{geom}}|$.
Take the inverse image of the cyclic order-thirty-one subgroup in
the boundary monodromy group.  A Sylow thirty-one subgroup of this
inverse image maps onto the cyclic quotient.  There are no elements
of order $31^2$ in $\mathrm{GL}_{10}(2)$, since the multiplicative
order of $2$ modulo $31^2$ is 155.  Every nonidentity element of this
Sylow subgroup therefore has order thirty-one; choose one with
nonidentity quotient image and denote it by $g$.
It acts trivially on the graph submodule, since $31\nmid|S_5|$.
Its characteristic polynomial consequently has the form
\begin{equation}\label{eq:sextic-order31}
 \chi_g(T)=(T-1)^5 f(T),
\end{equation}
where $f$ is an irreducible polynomial of degree five whose roots
are primitive thirty-first roots of unity.
This polynomial is not reciprocal: inversion does not belong to
the orbit of squaring on primitive thirty-first roots, since
$\langle2\rangle\subset(\mathbf Z/31)^\times$ has odd order five.
As before, the specialization is inside the full-rank Lang locus,
so these are elements of the same generic monodromy image.

Apply Hering's classification of transitive linear groups in the
form of Churikov--Vasil'ev--Zvezdina (2022), Lemma~3.3, arXiv p.~6,
which was opened and checked.  In dimension ten over $\mathbf F_2$,
its alternatives, unless the group contains $\mathrm{SL}_{10}(2)$,
place the group in one of
\[
 N(\mathrm{Sp}_{10}(2)),\qquad
 \Gamma\mathrm L_5(4),\qquad
 \Gamma\mathrm L_2(32),\qquad
 \Gamma\mathrm L_1(1024).
\]
The exceptional alternatives have dimensions four or six or odd
characteristic and do not occur here; the $G_2$ alternative has
dimension divisible by six.
The symplectic normalizer preserves its unique invariant alternating
form up to a scalar, and over $\mathbf F_2$ the scalar is one.
Every one of its characteristic polynomials is reciprocal, contrary
to \eqref{eq:sextic-order31}.
An element of order thirty-one in $\Gamma\mathrm L_5(4)$ lies in
$\mathrm{GL}_5(4)$, so its fixed space has even dimension over
$\mathbf F_2$, contrary to the fixed dimension five in
\eqref{eq:sextic-order31}.
Finally, neither $\Gamma\mathrm L_2(32)$ nor
$\Gamma\mathrm L_1(1024)$ has order divisible by seven, which has
already been forced by the quintic quotient.
All proper alternatives are excluded.  Therefore
$G_{\mathrm{geom}}=\mathrm{SL}_{10}(2)=\mathrm{GL}_{10}(2)$.
The arithmetic image contains this group and is contained in it,
so it is the same group.
\end{proof}
